%% Elevark fra mattelab.no – laget automatisk av bygg.py. Rediger gjerne til eget bruk.
%% Kompiler med lualatex.
\documentclass[11pt,a4paper]{article}
\usepackage[norsk]{babel}
\usepackage[a4paper,margin=2.2cm,top=2cm,bottom=2.4cm]{geometry}
\usepackage{fontspec}
\usepackage{unicode-math}
\setmainfont{TeX Gyre Pagella}
\setsansfont{TeX Gyre Heros}[Scale=MatchLowercase]
\setmathfont{TeX Gyre Pagella Math}
\usepackage{xcolor}
\definecolor{blekk}{HTML}{16213A}
\definecolor{dempet}{HTML}{566077}
\definecolor{rute}{HTML}{CFC9B8}
\definecolor{m1}{HTML}{2F4AB0}
\definecolor{m2}{HTML}{99266D}
\usepackage[most]{tcolorbox}
\usepackage{qrcode}
\usepackage{needspace}
\usepackage{fancyhdr}
\usepackage[hidelinks]{hyperref}
\pdfvariable suppressoptionalinfo \numexpr 1+2+4+8+16+32+64+128+256+512\relax
\setlength{\parindent}{0pt}
\setlength{\parskip}{0.4em}
\pagestyle{fancy}
\fancyhf{}
\renewcommand{\headrulewidth}{0pt}
\fancyfoot[L]{\sffamily\footnotesize\color{dempet}mattelab.no · Eksponentiell og logistisk vekst · elevark uten fasit}
\fancyfoot[R]{\sffamily\footnotesize\color{dempet}Side \thepage}

\newcommand{\fase}[2]{\tcbox[on line, colframe=#2, colback=white, coltext=#2, boxrule=0.6pt, arc=2pt,
  boxsep=0pt, left=3pt, right=3pt, top=1.5pt, bottom=1.5pt, fontupper=\sffamily\bfseries\footnotesize]{#1}}
\newcommand{\steg}[4]{%
  \par\noindent{\color{rute}\rule{\linewidth}{0.4pt}}\par\nopagebreak
  \noindent\begin{minipage}[t]{2.1cm}\vspace{0pt}\fase{#1}{#2}\end{minipage}%
  \begin{minipage}[t]{\dimexpr\linewidth-2.1cm\relax}\vspace{0pt}#3\end{minipage}\par
  \vspace{#4}}
\newcommand{\del}[4]{%
  \par\noindent{\color{rute}\rule{\linewidth}{0.4pt}}\par\nopagebreak
  \noindent\begin{minipage}[t]{1.1cm}\vspace{0pt}\fase{#1}{#2}\end{minipage}%
  \begin{minipage}[t]{\dimexpr\linewidth-1.1cm\relax}\vspace{0pt}#3\end{minipage}\par
  \vspace{#4}}
\newcommand{\oppgave}[2]{\par\Needspace{12\baselineskip}\vspace{0.6em}%
  {\fontsize{14}{17}\selectfont\bfseries{\color{m1}#1}\enspace #2}\par\vspace{0.2em}}

\begin{document}
\begin{minipage}[t]{0.72\linewidth}\vspace{0pt}
{\sffamily\small\color{dempet}Elevark · 1T og R1 · mattelab.no}\par\vspace{0.3em}
{\fontsize{24}{28}\selectfont\bfseries Eksponentiell og logistisk vekst}\par\vspace{1em}
{\sffamily\small Navn: \rule{4.6cm}{0.4pt}\quad Klasse: \rule{1.5cm}{0.4pt}\quad Dato: \rule{1.9cm}{0.4pt}}
\end{minipage}\hfill
\begin{minipage}[t]{0.22\linewidth}\vspace{0pt}\raggedleft
\qrcode[height=2.3cm]{https://mattelab.no/utforskere/vekst.html}\par\vspace{0.2em}
{\sffamily\scriptsize\color{dempet}Utforskeren}
\end{minipage}

\vspace{1em}
{\sffamily\small\color{dempet}Utforskeren ligger på \url{https://mattelab.no/utforskere/vekst.html}}\par
Hver oppgave har tre steg. \textbf{Gjett} først, uten å bruke utforskeren. \textbf{Test} så i utforskeren. \textbf{Forklar} til slutt med egne ord hvorfor det blir slik.

\oppgave{1}{Penger i banken}
\steg{Gjett}{m1}{Du setter inn $1000$ kr på en konto med $5\,\%$ rente per år. Hvor mye tror du står på kontoen etter $10$ år? Hvor lang tid tar det før beløpet er doblet?}{2.2cm}
\steg{Test}{m2}{Velg eksponentiell modell med $N_0 = 1000$, $p = 5$ og $t = 10$. Dra $t$ for å finne når beløpet passerer $2000$ kr.}{2.6cm}
\steg{Forklar}{blekk}{Hvorfor blir det mer enn $1000 + 10 \cdot 50 = 1500$ kr? Regn ut doblingstiden med formelen.}{3.4cm}

\oppgave{2}{Doblingstid og startverdi}
\steg{Gjett}{m1}{En bakteriekultur vokser med $20\,\%$ per time. Tror du doblingstiden blir kortere, lengre eller lik om vi starter med $50$ bakterier i stedet for $10$?}{2.2cm}
\steg{Test}{m2}{Velg eksponentiell modell med $N_0 = 10$ og $p = 20$. Les av doblingstiden. Endre så $N_0$ til $50$. Prøv også $p = 5$ og $p = 100$.}{2.6cm}
\steg{Forklar}{blekk}{Vis med formelen at doblingstiden ikke avhenger av $N_0$. Sammenlign doblingstidene med tommelfingerregelen $\frac{70}{p}$. Når er regelen god?}{3.4cm}

\oppgave{3}{Medisin i blodet}
\steg{Gjett}{m1}{En pasient får $400$ mg av en medisin. Mengden i blodet minker med $15\,\%$ per time. Hvor mye tror du er igjen etter $6$ timer? Hva er halveringstiden?}{2.2cm}
\steg{Test}{m2}{Velg eksponentiell modell med $N_0 = 400$, $p = -15$ og $t = 6$.}{2.6cm}
\steg{Forklar}{blekk}{Hvorfor er vekstfaktoren $0{,}85$? Hvorfor blir mengden aldri helt null i denne modellen?}{3.4cm}

\oppgave{4}{Når bremser veksten?}
\steg{Gjett}{m1}{En dam har plass til omtrent $200$ fisker. Vi setter ut $10$ fisker, og bestanden vokser først med omtrent $20\,\%$ per år. Hvor mange fisker tror du det er etter $20$ år? Når vokser bestanden raskest?}{2.2cm}
\steg{Test}{m2}{Vis begge modellene med $N_0 = 10$, $p = 20$, $K = 200$ og $t = 20$. Sammenlign modellene ved $t = 5$, $t = 10$ og $t = 20$. Finn vendepunktet.}{2.6cm}
\steg{Forklar}{blekk}{Hvorfor er modellene nesten like i starten? Hvilken modell er mest realistisk etter $20$ år, og hvorfor?}{3.4cm}

\oppgave{5}{For mange fra start}
\steg{Gjett}{m1}{Hva tror du skjer i den logistiske modellen hvis vi setter ut $300$ fisker i en dam med bæreevne $200$?}{2.2cm}
\steg{Test}{m2}{Velg logistisk modell med $N_0 = 300$, $p = 20$ og $K = 200$.}{2.6cm}
\steg{Forklar}{blekk}{Se på formelen for $N(t)$. Hvilket fortegn har $c$ nå? Forklar ut fra formelen hvorfor $N(t)$ minker mot $K$.}{3.4cm}

\end{document}
